% !TeX root = ../main.tex

\chapter{总结与展望}
\label{ch:conclusion}

\section{工作总结}

本课题围绕遥感场景下的微小船舶检测，完成了研究综述、多模型实验对比和系统实现三条主线工作。研究层面，对通用目标检测、小目标检测和遥感船舶检测相关方法进行了梳理，并在统一评测协议下组织了 DRENet、FCOS、YOLO26 的对比实验。工程层面，基于 React、FastAPI 和 SQLite 构建了可运行的检测系统，打通了任务提交、异步推理、结果可视化与历史回看链路。

从结果上看，DRENet 与 YOLO26 在 AP50 上接近，但特征不同：DRENet 召回更高，适合漏检敏感场景；YOLO26 在 AP50-95、Precision 与 F1 上更均衡，更适合作为系统默认候选。FCOS 整体表现略弱于前两者，但保留了 anchor-free 路线的对照价值，补充了不同检测范式的比较维度。

从系统实现上看，本课题没有把系统部分仅作为界面展示，而是把它作为实验结果管理与算法能力封装的平台。通过前后端分层、异步任务机制和持久化设计，系统已经能够支持真实模型联调、任务回放和答辩演示。

整体上，课题已形成“统一实验协议 + 真实结果支撑 + 系统闭环落地”的完整成果链路。工作重点不在提出全新算法，而在于给出可复核、可对比、可演示的完整实现路径。

\section{不足分析}

当前工作仍存在几个边界需要客观说明。

第一，输入尺寸敏感性分析目前覆盖 YOLO26 与 FCOS，DRENet 在本地插件路径下使用非 512 输入会触发 shape mismatch，因此未纳入统一尺寸主表。第二，FCOS 的正式结果来自 MMDetection 默认输入策略，训练与效率口径与 DRENet、YOLO26 并非严格同尺度一致，更适合承担参考模型角色。第三，系统已具备完整运行链路，但插件化部署规范与跨机器迁移说明仍可进一步完善，部分模型接入仍依赖本地代码路径。第四，融合模式已可稳定运行，但尚未形成统一离线 AP 主表。

这些问题属于可继续优化的工程和实验边界，不影响当前主线结论成立。

\section{未来展望}

后续工作可沿以下方向推进：

其一，补齐融合模式的统一离线评测结果，建立单模型对比与融合展示之间的量化联系。其二，围绕 DRENet 与 YOLO26 的误检-漏检取舍，继续优化阈值、后处理与融合策略。其三，进一步统一 FCOS 的输入尺度与训练策略，提升跨模型比较的一致性。其四，完善模型插件化接入和部署规范，降低跨机器迁移成本。其五，在更多遥感场景与数据集上验证方法和系统的泛化能力。

\section{全文结论}

课题完成了面向遥感微小船舶检测的实验与系统一体化实现。实验结果表明，在 LEVIR-Ship 与统一评测协议下，DRENet 更偏召回优先，YOLO26 更偏精度与部署均衡，FCOS 提供了 anchor-free 路线对照。系统结果表明，将算法能力服务化并任务化管理后，能够显著提升结果展示、回放和工程交付的完整性。现有实验记录、分析框架与系统实现，为后续性能优化、评测补齐和部署规范化提供了可继续迭代的基础。
